\section{Computability}
\label{ch:computability}

This chapter collects the computability-theoretic ingredients of the proof: classical
statements with no quantum content, several of which partially exist in Mathlib. Like
the results of Section~\ref{ch:required-results}, these are external results that the
project may initially assume as axioms, and each entry carries the same
formalization-effort estimate. The concrete machine model and its canonical codes come
first; then the toolkit of resource-bounded computability results; the recursive
compression lemma, which builds on the toolkit, isolates the self-referential part of
the proof.

\subsection{The machine model}
\label{sec:machine-model}

The machine-level side of the project is formalized directly: a concrete multi-tape
machine model with several input tapes, a canonical first-order syntax of machine
\emph{codes} with an exact binary serialization, and a reference evaluator. 

\begin{definition}[Multi-input Turing machine]
\label{def:multiinput-tm}
\lean{Turing.MultiInputTM, Turing.MultiInputTM.ComputesInTimeAndSpace}
\leanok
A \emph{multi-input Turing machine} with $i$ input tapes and $w$ work tapes has $i$
two-way read-only input tapes (heads clamped one blank cell beyond either end of the
input), $w$ bidirectional work tapes over an alphabet extended with a blank symbol, and
a write-only output stream. A transition may write and move on every work tape, move
every input head, emit at most one output symbol, and either continue or halt. Time is
counted in steps, and space as the number of work-tape cells visited. 

At $i = 1$ this
is exactly the CSLib multi-tape model.
\end{definition}

\begin{definition}[Machine codes]
\label{def:machine-code}
\uses{def:multiinput-tm}
\lean{Turing.Code, Turing.Code.toTM, Turing.encodeCode, Turing.decodeCode,
Turing.codeSize}
\leanok
A \emph{machine code} is first-order data: a header (number of work tapes, alphabet
size, number of states, start state) and a dense transition table with one entry per
observation, in a fixed canonical order. A well-formed code is interpreted as a
multi-input machine, with binary input and output through two reserved alphabet
symbols. Codes serialize exactly: the \emph{description} $\alpha \in \{0,1\}^*$ of a
code has size $\abs{\alpha}$, and \emph{total decoding} interprets every string ---
$[\alpha]_i$ denotes the coded machine with $i$ input tapes described by $\alpha$, with
every malformed string denoting one fixed machine that immediately rejects.
\end{definition}

\begin{definition}[Runs of coded machines]
\label{def:code-run}
\uses{def:machine-code}
\lean{Turing.Code.Produces, Turing.Code.evalWithin}
\leanok
Let $c$ be a machine code with $i$ input tapes, $x = (x_1, \ldots, x_i)$ bit strings and $y \in
\{0,1\}^*$. The code $c$ \emph{produces} $y$ on $x$ if its run on $x$ halts with output
$y$ (in some time and space). The \emph{budgeted evaluation} of $c$ on $x$ with budget
$T$ returns the output if $c$ has halted within $T$ steps, and a timeout marker
otherwise; it is total and executable, and the two notions agree: $c$ produces $y$ if
and only if some budget suffices to observe it.
\end{definition}

\begin{lemma}[Canonical serialization]
\label{lem:code-canonical}
\uses{def:machine-code}
\lean{Turing.decodeCodeExact_encodeCode, Turing.encodeCode_injective,
Turing.decodeCodeExact_sound}
\leanok
Exact decoding inverts encoding, distinct codes have distinct descriptions, and every
description accepted by the exact decoder is the canonical encoding of the code it
parses to.
\begin{proof}
\leanok
Compositional, from paired prefix and soundness lemmas for each field of the
description format.
\end{proof}
\end{lemma}

The following two statements are the machine-level counterparts of
Lemma~\ref{lem:universal-tm}. 

\begin{lemma}[Universal machine]
\label{lem:universal-machine}
\uses{def:machine-code, def:code-run}
\lean{Turing.exists_universalCode}
\leanok
For each arity $i$ there are a machine code $U$ with $i + 1$ input tapes and a
polynomial $P$ such that for all $\alpha \in \{0,1\}^*$, all inputs $x$ and all $y$:
$U$ on $(\alpha; x)$ produces $y$ if and only if $[\alpha]_i$ on $x$ produces $y$; and
whenever $[\alpha]_i$ halts on $x$ within $t$ steps, $U$ halts on $(\alpha; x)$ within
$P(\abs{\alpha} + t)$ time and space.
\end{lemma}

\begin{proof}
Not planned: issue \#17 was closed on 2026-09-24, and the Lean statement is left as a
\texttt{sorry}. The main theorem does not use this lemma; its universal machine is the ambient
one of Lemma~\ref{lem:universal-tm}.
\end{proof}

\begin{lemma}[Bounded universal machine]
\label{lem:bounded-universal-machine}
\uses{def:machine-code, def:code-run}
\lean{Turing.exists_boundedUniversalCode}
\leanok
For each arity $i$ there are a machine code $U$ with $i + 2$ input tapes and a
polynomial $P$ such that on input $(\alpha;\, T \text{ in binary};\, x)$, the machine
$U$ always halts within $P(\abs{\alpha} + T)$ steps, outputting the canonically encoded
budgeted evaluation of $[\alpha]_i$ on $x$ with budget $T$.
\end{lemma}

\begin{proof}
Not planned: issue \#18 was closed on 2026-09-24, and the Lean statement is left as a
\texttt{sorry}. The main theorem does not use this lemma; its universal machine is the ambient
one of Lemma~\ref{lem:universal-tm}.
\end{proof}

\subsection{Efficient computability toolkit}
\label{sec:rr-computability}

\effortHard{} (infrastructure: requires a cost model for computation)

\begin{lemma}[Efficient universal Turing machine]
\label{lem:universal-tm}
\lean{MIPRE.Cost.exists_efficient_universal, MIPRE.Cost.exists_clocked_universal}
\leanok
There is a program $u$ and a polynomial $p$ such that $\phi_u(e, x) = \phi_e(x)$ for
all $e, x$, with $\mathrm{runtime}(u, (e,x)) \leq
p\bigl(\abs{e}, \abs{x}, \mathrm{runtime}(e, x)\bigr)$.
\end{lemma}

\begin{proof}
\leanok
Both the universal program and its clocked variant are self-interpreters of the ambient
list language, written in the language itself (\texttt{MIPRE.Cost.Machine.univProg},
\texttt{MIPRE.Cost.Machine.univTProg}). Evaluation is first presented as an abstract
machine on configurations (a control ``evaluate $p$'' or ``return $v$'', an environment
and a stack of frames) whose steps are charged exactly the costs of the evaluation rules,
so that a derivation of cost $t$ corresponds to a machine run of at most $3t$ steps
reaching a final configuration, and conversely every run reaching a final configuration
comes from a derivation. Configurations are encoded as data and the step function is
implemented by a fixed program (\texttt{stepProg}) at cost quadratic in the size of the
configuration. Along the run simulating a derivation of cost $t$ on the input $x$, every
configuration has size polynomial in $\abs{e} + \abs{x} + t$: values are results of
sub-derivations, hence of size at most $t$; environments and stacks grow by at most one
entry per unit of cost. The universal program iterates \texttt{stepProg} until a final
configuration is reached, so its runs are polynomially bounded, and any of its halting
runs is a machine run reaching a final configuration, hence a derivation with the same
result. The clocked variant additionally carries, in unary, the remaining cost budget
(decremented by the exact cost of each step), a step budget $3k+1$ guaranteeing
termination, and a size guard $\Theta$, a fixed polynomial in $k + \abs{e} + \abs{x}$
dominating the configuration sizes of any run of cost at most $k$; it returns the timeout
value when a budget or the guard is exceeded. The guard is checked on the configuration
produced by each step, whose size is bounded by the cost of producing it, so the loop
only handles configurations of polynomial size even when the simulated run does not halt.
\end{proof}

\begin{lemma}[Efficient s-m-n]
\label{lem:smn}
\lean{MIPRE.Cost.hardcode_time, MIPRE.Cost.hardcode_time_rev, MIPRE.Cost.hardcode_size, MIPRE.Cost.smn_polyTime}
\leanok
There is a polynomial-time computable $s$ such that $\phi_{s(e,x)}(y) = \phi_e(x, y)$
for all $e, x, y$, with the runtimes of $e$ and $s(e,x)$ polynomially equivalent.
\end{lemma}

\begin{proof}
\leanok
In the ambient list language, $s(e, x)$ is the program that pairs the literal $x$ with
its input and runs $e$ on the pair (\texttt{MIPRE.Cost.hardcode}). The time overhead is
additive and linear in $\abs{x}$ and the input size, in both directions, and the
description of $s(e, x)$ is that of $e$ plus $\abs{x}$ plus a constant. The map $s$
itself is computed in linear time by a fixed program assembling the description
(\texttt{smnProg}).
\end{proof}

\begin{lemma}[Efficient Kleene recursion theorem]
\label{lem:kleene}
\uses{lem:universal-tm,lem:smn}
\lean{MIPRE.Cost.efficient_fixed_point}
\leanok
For every polynomial-time computable $f : \{0,1\}^* \to \{0,1\}^*$ there is a program
$e$ with $\phi_e = \phi_{f(e)}$ and the runtime of $e$ polynomially bounded in terms
of that of $f(e)$. (This is the direction used by Lemma~\ref{lem:recursive-compression};
\cite{MNY} states the runtimes as polynomially equivalent, and the converse direction
would require a simulator that is never faster than the simulated program.)
\end{lemma}

\begin{proof}
\leanok
The classical construction. Let $G$ be the program that on input $(x, v)$ computes
$s(x, x)$ (Lemma~\ref{lem:smn}), applies $f$ to it, and runs the resulting program on $v$
through the universal machine (Lemma~\ref{lem:universal-tm}); put $e = s(G, G)$. Then $e$
on $v$ is $G$ on $(G, v)$, which runs $f(s(G, G)) = f(e)$ on $v$: the two programs halt
on the same inputs with the same results (the backward direction inverts the run of $G$
step by step and uses that the universal machine halts only if the simulated program
does), and a run of $f(e)$ of time $t$ gives a run of $e$ within $p(\abs{v}, t)$, where
$p$ is the universal machine's overhead polynomial shifted by the fixed size of $f(e)$,
plus a linear term for the copies of the input and a constant for the fixed computation
of $f(e)$.
\end{proof}

\paragraph{Source.} Folklore; stated in this efficient form in~\cite{MNY}, Lemmas
2.1--2.3.

\paragraph{Comments.} Mathlib contains the non-resource-bounded versions (universal
machine, s-m-n and the Rogers fixed-point theorem for \texttt{Nat.Partrec.Code}). The
efficient versions require a cost model for the chosen machine model; this is real
infrastructure work and a prerequisite for everything involving ``polynomial-time''.
The formalization's ambient model is not a Turing machine but a small first-order
list-processing language (\texttt{MIPRE.Cost.Prog}, over binary trees) with a unit-cost
semantics in which reading a value costs its size; it is polynomially equivalent to
Turing-machine time, programs are their own tree serializations (so ``a program $e$''
and ``$\abs{e}$'' are literal), and the s-m-n theorem is a linear-overhead syntactic
operation. Mathlib's \texttt{Turing.ToPartrec.Code} was tried first and rejected: its
programs cannot rebuild a list in polynomial time.
A pragmatic first pass can replace ``polynomial time'' by ``computable'' wherever
Section~\ref{sec:departures} permits, deferring the cost model. The machine-level
counterpart, for the concrete model of Section~\ref{sec:machine-model}, is stated as
Lemmas~\ref{lem:universal-machine} and~\ref{lem:bounded-universal-machine}; the
pipeline consumes only the present ambient version.

\subsection{The abstract compression lemmas}
\label{sec:rr-compression}

\effortEasy{} (given the toolkit of Section~\ref{sec:rr-computability})

\begin{lemma}[Recursive compression]
\label{lem:recursive-compression}
\uses{def:succinct,lem:universal-tm,lem:smn,lem:kleene}
\lean{MIPRE.Cost.recursive_compression}
\leanok
Let $f : \{0,1\}^* \to \N \cup \{\infty\}$ be a function and $A \subseteq \{x : f(x) <
\infty\}$ a nonempty language. Suppose there is a polynomial-time computable function
$\compr : \{0,1\}^* \times \N \to \{0,1\}^*$ such that whenever $\hat{x}$ is a succinct
description (Definition~\ref{def:succinct}) of a string $x$ and $y = \compr(\hat{x},
n)$:
\begin{enumerate}
\item $f(y) \geq \max\{ f(x), n \}$, and
\item if $x \in A$ then $y \in A$.
\end{enumerate}
Then there is a polynomial-time computable function $g$ reducing the halting problem to
$A$: for every Turing machine $e$, if $\phi_e$ halts on the empty input then $g(e) \in
A$, and otherwise $f(g(e)) = \infty$.
\end{lemma}

\begin{proof}
\leanok
The construction of~\cite{MNY}, with two harmless modifications: the ``next level'' of
the recursion is $2n + 1$ instead of $n + 1$ (prepending a bit to a binary numeral is one
step in the list language, an increment is not), and the threshold below is explicit.
Let $\beta$ be the fixed \emph{bit-query} program that, on input $(((c', e), n), m)$,
runs the universal machine on $c'$ with input $(e, n)$ to obtain a string $y$ and outputs
the $m$-th bit of $y$ (or $2$ if $m \geq \abs{y}$). Let $a$ be the polynomial-time
\emph{decider} that on $(c', (e, n))$ outputs a fixed $y_0 \in A$ if $e$ halts on the
empty input within $\log n$ steps, and $\compr\bigl((s(\beta, ((c', e), 2n+1)), n),
n\bigr)$ otherwise, where $s$ is the s-m-n map of Lemma~\ref{lem:smn}. By
Lemma~\ref{lem:kleene} there is a program $c$ with $\phi_c(e, n) = \phi_a(c, (e, n))
=: h(e, n)$, running in time polynomial in $\abs{e} + \log n$. Hence for $n$ above a
threshold $r(e) = 2^{K + 1 + \abs{e}}$ (with $K$ a constant depending only on the
polynomial overheads), $(s(\beta, ((c, e), 2n+1)), n)$ is a succinct description of
$h(e, 2n+1)$: the string has length at most $2^n$ and the bit-query program answers within
the budget of Definition~\ref{def:succinct}. The reduction is $g(e) = h(e, r(e))$. If $e$
does not halt, then for every $n \geq r(e)$ the decider compresses, so $f(h(e, n)) \geq
\max\{f(h(e, 2n+1)), n\}$, and iterating gives $f(h(e, n)) \geq n + k$ for all $k$,
i.e.\ $f(g(e)) = \infty$. If $e$ halts in $T$ steps, then $h(e, n) = y_0 \in A$ once
$n \geq 2^T$, and by downward induction along the levels $n, 2n+1, 4n + 3, \ldots$ the
second property of $\compr$ gives $h(e, n) \in A$ for every $n \geq r(e)$, in particular
$g(e) \in A$.
\end{proof}

\paragraph{Source.} Marks, Nezhadi and Yuen~\cite{MNY}, Lemma 5.1 (the variant of
their Lemma 3.1 with compression parameter $n$); the technique is implicit
in~\cite{DRS08,FJVY19,JNVWY20}. The proof is a fixed-point construction using the
efficient universal machine, s-m-n and Kleene recursion theorems
(Section~\ref{sec:rr-computability}).

\paragraph{Comments.} This is the entire ``self-referential'' part of the proof,
isolated in a statement with no quantum content. In the application
(Section~\ref{sec:ps-halting}), $A$ is the set of normal form verifier descriptions
with $\synval = 1$ and finite entanglement requirement, and $f(\verifier) =
\Ent(\verifier, \frac12)$ (suitably formalized); the compression procedure is
Theorem~\ref{thm:compression}. The value-form pipeline of
Section~\ref{ch:proof-structure} uses Lemma~\ref{lem:compressible-criterion} below
instead; this lemma is retained as the entanglement-form route
(Remark~\ref{rem:entanglement-form}), where the factor $\frac12$ in the entanglement
bound of compression requires the variant with hypothesis $f(y) \geq \frac12
\max\{f(x), G(n)\}$ for a function $G$ growing faster than the levels. Note that~\cite{MNY} also proves a converse (their
Lemma 4.1): a halting reduction of this kind exists \emph{only if} a compression
procedure does, which is evidence that this decomposition of the proof is canonical.
The lemma and its proof are a self-contained computability project, an ideal early
milestone.

\begin{lemma}[Compressibility criterion]
\label{lem:compressible-criterion}
\uses{def:succinct,lem:universal-tm,lem:smn,lem:kleene}
\lean{MIPRE.Cost.compressibility_criterion, MIPRE.Cost.compressibility_criterion_levels}
\leanok
Let $A, B \subseteq \{0,1\}^*$ be languages with distinguished elements
$y_{\mathrm{yes}} \in A$ and $y_{\mathrm{no}} \in B$, and suppose that the complement of
$B$ is recursively enumerable: there is a program $S$ that halts on input $x$ if and only
if $x \notin B$. Suppose there is a polynomial-time computable function $\compr : \{0,1\}^*
\times \N \to \{0,1\}^*$ such that whenever $\hat{x}$ is a succinct description
(Definition~\ref{def:succinct}) of a string $x$ with $\abs{x} \leq n + 1$ and
$y = \compr(\hat{x}, n)$:
\begin{enumerate}
\item if $x \in A$ then $y \in A$, and
\item if $x \in B$ then $y \in B$.
\end{enumerate}
Then there is a polynomial-time computable function $g$ reducing the halting problem to
$(A, B)$: for every Turing machine $e$, if $\phi_e$ halts on the empty input then $g(e)
\in A$, and otherwise $g(e) \in B$.
\end{lemma}

\begin{proof}
\leanok
The construction of Lemma~\ref{lem:recursive-compression}, with a third branch. Let
$\beta$ be the bit-query program of that proof, and let $\sigma$ be the program that, on
$((c', (e, (m, m))), \cdot)$, runs the universal machine on $c'$ with input $(e, (m, m))$
to obtain a string $x$ and then runs $S$ on $x$. Let $a$ be the polynomial-time decider
that on $(c', (e, (m, n)))$ outputs $y_{\mathrm{yes}}$ if $e$ halts on the empty input
within $\log n$ steps; otherwise $y_{\mathrm{no}}$ if $s(\sigma, (c', (e, (m, m))))$
halts within $\log n$ steps; and otherwise $\compr\bigl((s(\beta, ((c', e), (m,
2n+1))), n), n\bigr)$, where $s$ is the s-m-n map of Lemma~\ref{lem:smn}. By
Lemma~\ref{lem:kleene} there is a program $c$ with $\phi_c(e, (m, n)) = \phi_a(c, (e, (m,
n))) =: h(e, m, n)$, running in time polynomial in $\abs{e} + \log m + \log n$. As
before, for $n$ above a threshold $r(e) = 2^{K + 1 + \abs{e}}$ and $m = r(e)$, the pair
$(s(\beta, ((c, e), (m, 2n+1))), n)$ is a succinct description of $h(e, m, 2n+1)$; the
start level $m$ is carried as data so that the threshold can be chosen from the running
time of $a$ afterwards. The reduction is $g(e) = h(e, r(e), r(e)) =: x_e$; the second
branch tests, at each level it is run at, whether $S$ halts on $x_e$. Two facts are used
implicitly below and hold along the chain: $\abs{m} \leq \abs{n}$ at levels $n \geq r(e)
\geq m$, which is what keeps $K$ out of its own defining polynomial; and ``$e$ halts within
$\log n$ steps'' is monotone in $n$, so the first branch, once it fires, fires at every
higher level.

If $e$ does not halt, suppose $x_e \notin B$. Then $S$ halts on $x_e$ in some time $T$,
so at every level $n$ with $\log n \geq T$ the decider outputs $y_{\mathrm{no}} \in B$,
and by downward induction along the levels $r(e), 2r(e)+1, 4r(e)+3, \ldots$ from such a
level, property 2 of $\compr$ gives $h(e, r(e), n) \in B$ at every level of the chain, in particular
$x_e \in B$, a contradiction. If $e$ halts in $T$ steps, then $h(e, r(e), n) =
y_{\mathrm{yes}} \in A$ once $\log n \geq T$, and by downward induction property 1 gives
$h(e, r(e), n) \in A$ at every level of the chain, provided the second branch never fires at a level
where the first has not. If it fired at such a level $n_1$, then $S$ halts on $x_e$, so
$x_e \notin B$; but the output $y_{\mathrm{no}} \in B$ at level $n_1$ and property 2, by
downward induction along the levels between $r(e)$ and $n_1$ (at which the first branch
does not fire either), give $x_e \in B$, a contradiction. Hence $g(e) = x_e \in A$.
\end{proof}

\paragraph{Source.} Lin~\cite{Lin25}, the compression criterion for $\RE$-complete
problems (direction ``weakly compressible implies $\RE$-complete''), whose
self-referential sequence has exactly these three branches; the formulation through
succinct descriptions is that of~\cite{MNY} and of Lemma~\ref{lem:recursive-compression}.
Lin's theorem is an equivalence: given that the no-instances form a $\coRE$ set, a
decision problem is $\RE$-complete if and only if it is compressible.

\paragraph{Comments.} This is the lemma the main theorem uses
(Section~\ref{sec:ps-halting}): $A$ is the set of normal form verifier descriptions with a
value-$1$ PCC strategy, $B$ the set of those with $\valstar \leq \frac12$, and $S$ the
$\valstar$ half of the enumeration of Lemma~\ref{lem:value-lower-approx}. Compared with
Lemma~\ref{lem:recursive-compression}, the growth hypothesis on a measure is replaced by
the semidecidability of the complement of $B$ --- both directions of the specification of
$S$ are used --- and nothing about entanglement is required of the compression procedure.
The Lean proof reuses the toolkit of Section~\ref{sec:rr-computability}: the search
branch runs the semidecision procedure through the clocked universal machine, and the
start level of the recursion is carried as data next to the current level.

The Lean file proves a \emph{per-level} form first,
\texttt{compressibility\_criterion\_levels}, and derives the statement above from it. There
the classes are families $A_n, B_n$ indexed by the level $n$ of the recursion, the
semidecision procedure receives the pair $(x, n)$, and the compression hypothesis reads: for
$n \geq n_0$, if $\hat x$ describes $x$ succinctly with parameter $n$ and $2\abs{\hat x}
\leq n$, then $x \in A_{2n+1}$ implies $\compr(\hat x, n) \in A_n$, and likewise for $B$;
the conclusion places $g(e)$ in $A_{r(e)}$ or $B_{r(e)}$ at the explicit start level $r(e)
= 2^{K + 1 + \abs{e}}$, which is at least $n_0$. This is exactly what the proof uses, and
it is the form an instantiation by gap-preserving compression can meet
(Remark~\ref{rem:compression-abstract}): $\compr$ runs in time polynomial in the
\emph{length} of $(\hat x, n)$, hence polylogarithmic in $n$, while a single bit query to
$\hat x$ costs up to $n$, so $\compr$ cannot read the string $x$ it compresses and can only
embed $\hat x$ into its output, for a reader whose budget is polynomial in $n$ --- the
decider of a game at index $n$. The size slack $2\abs{\hat x} \leq n$ keeps that output
short, and the levels are what let the class of the output be judged at the index where
it is played.

The hypothesis $\abs{x} \leq n + 1$ is not an extra demand on $\compr$ but a \emph{weakening} of
what it must cover, and it was added because obligation O4 cannot be met without it. The
proof already has it: the described string is the output of a run whose cost the threshold
constant bounds by $n + 1$, so $\abs{x} \leq \abs{\mathrm{encode}(x)} \leq t \leq n + 1$,
and until this was carried over the proof discarded it into
Definition~\ref{def:succinct}'s $\abs{x} \leq 2^n$ through $n + 1 \leq 2^n$. Why O4 needs
it: the class $A$ of the halting reduction asks for a value-$1$ PCC strategy at the answer
bound $\poly(n, \lambda_x)$ read off the described string's own compression parameter
$\lambda_x$, while the only source of such a strategy for a compressed verifier is the
completeness clause of Theorem~\ref{thm:compression}, which takes its input at the bound
$(2^n)^\lambda$ with $\lambda$ the parameter the compressor \emph{writes into its output} and so
bounded by the compressor's own output size. Nothing in Definition~\ref{def:lambda-bounded}
constrains $\lambda_x$ --- for a compression procedure whose compressed sampler does not vary
with $\lambda$, membership in either class is insensitive to it --- so with only
$\abs{x} \leq 2^n$ the two bounds cannot be related by any single choice of $\lambda$. With
$\abs{x} \leq n + 1$ we get $\lambda_x < 2^{n+1}$, hence an answer bound of
$2^{O(n)}$ against $(2^n)^\lambda$, which $\lambda \geq 2\deg(\mathrm{bound}) + 1$ already
dominates for $n$ above a threshold read off $\mathrm{bound}$ as well; the Lean says so and
checks it (\texttt{MIPRE.Halting.exists\_ansBound\_le}).

That settles the class $A$ direction. The class $B$ direction needs the same inequality read
the other way --- the soundness clause of Theorem~\ref{thm:compression} wants the value
bounded at $(2^n)^\lambda$, and enlarging the answer bound can only raise a value --- so a
single $\lambda$ serves both only once the described string's own decider is known to reject
answers longer than its answer bound, which is what $\mathrm{poly}(n, \lambda_x)$ was chosen
to be. The paper's $\vhalt$ has that by construction --- Step~6 of its decider is an explicit
length check, spent in the amended proof of Lemma~\ref{lem:dhalt-values} --- and the
instantiation asks for it as a clause of membership in $A_n$ and $B_n$
(Remark~\ref{rem:compression-abstract}, item 1; \texttt{Verifier.RejectsLong}): the two
distinguished strings have it, the compressor's output has it by the length check of
Theorem~\ref{thm:compression}'s decider (\texttt{output\_rejects\_long}), and its violation is
enumerable (Lemma~\ref{lem:halting-semidecider}). The alternative, a well-formedness predicate
threaded through this lemma and instantiated to the same clause, would have left the classes
alone and touched this lemma's proof instead; the classes were chosen so that this lemma stays
as stated. Nothing more is asked of it.

\effortEasy

\begin{lemma}[Halting form of the two reduction lemmas]
\label{lem:halting-form}
\uses{lem:recursive-compression,lem:compressible-criterion}
\lean{MIPRE.Cost.recursive_compression_halting,
MIPRE.Cost.compressibility_criterion_halting}
\leanok
Lemma~\ref{lem:recursive-compression} and Lemma~\ref{lem:compressible-criterion} both
hold with the halting problem taken in the ambient formalization's own form: the reduction
may be presented as a computable map defined on the codes of a fixed G\"odel numbering of
the partial recursive functions, halting being halting of the code on input $0$, rather
than on the programs of the cost model in which the two lemmas are proved.
\end{lemma}

\begin{proof}
\leanok
\uses{lem:recursive-compression,lem:compressible-criterion}
Compose the reduction with a compiler from codes to programs of the cost model that is
computable on descriptions and preserves halting on the empty input: hardcode the code's
index into a fixed program for the universal partial function, which exists because that
partial function is partial recursive, and observe that the description of the result is a
fixed primitive recursive tree around the index. Only the interface changes; neither
lemma's content is touched, and the compiler carries no time bound, which is why this step
is available for the computable version of the main theorem and not for the
polynomial-time one.
\end{proof}

\paragraph{Comments.} This is the seam at which the development meets Mathlib's halting
problem, hence the form Theorem~\ref{thm:main} consumes: \texttt{Nat.Partrec.Code} is
what \texttt{Computable} is stated against, while the two reduction lemmas are proved
about programs of the cost model of Section~\ref{sec:rr-computability}, whose
polynomial-time predicate the fixed-point construction needs. Separating the two keeps the
cost model out of the main statement entirely.

\subsection{Deciding a cost budget, and semideciders on encoded values}
\label{sec:rr-cost-budget}

Lemma~\ref{lem:compressible-criterion} consumes semidecidability as an \emph{equivalence}:
its program $S$ halts on exactly the complement of $B$. When the criterion is instantiated
by the halting reduction that complement is a disjunction of two searches
(Remark~\ref{rem:compression-abstract}, item 3), one of them for an input on which a time
bound of Definition~\ref{def:lambda-bounded} fails. Recognizing such an input exactly ---
and not merely recognizing a program that runs for a long time --- is what the first lemma
here provides; the second moves the resulting semidecider off bit strings, the criterion's
$B$ being indexed by a level as well as a string.

\begin{lemma}[Halting within a cost budget is decidable]
\label{lem:cost-budget-decidable}
\lean{MIPRE.Cost.Machine.stepCostD, MIPRE.Cost.Machine.stepCostD_toData,
  MIPRE.Cost.Machine.costStep, MIPRE.Cost.Machine.costStep_iterate,
  MIPRE.Cost.Machine.evalForCostD, MIPRE.Cost.Machine.evalForCostD_eq_some,
  MIPRE.Cost.Machine.haltsWithin_of_evalForCostD,
  MIPRE.Cost.Machine.haltsWithin_iff_evalForCostD, MIPRE.Cost.Machine.haltsWithinB,
  MIPRE.Cost.Machine.primrec_haltsWithinB, MIPRE.Cost.decidableHaltsWithin}
\leanok
Whether a program of the cost model halts on a given input within a given cost is a
primitive recursive predicate of the program, the input and the cost; in particular it is
decidable.
\end{lemma}

\begin{proof}
\leanok
Carry the cost along the run. The cost the evaluation machine charges at a step is a
function of the configuration alone, so it is a function of the encoded configuration too
(\texttt{Machine.stepCostD}, dispatching on the same seven program tags and four frame tags
as the step function itself, and primitive recursive for the same reason); iterating the
step function paired with the running total, from the initial configuration and total $0$,
therefore computes the pair (configuration, accumulated cost) at every step
(\texttt{costStep\_iterate}). A derivation of cost $t$ is a machine run of at most $3t$
steps, so $3k + 1$ iterations suffice to reach the end of every run of cost at most $k$;
and once the end is reached nothing more is charged, a final configuration being fixed by
the step function and free. So the test is: iterate $3k + 1$ times, and report success if
the configuration reached is final \emph{and} the total is at most $k$
(\texttt{evalForCostD}). Completeness is the step count above; soundness is that the total
at a final configuration is the cost of the derivation that reached it, which is the
backward direction of the machine--derivation correspondence together with the freeness of
the steps after the end.
\end{proof}

\paragraph{Comments.} A step budget does not decide a cost budget, which is why this is not
a corollary of \texttt{Machine.runForD} (Remark~\ref{rem:compression-abstract}, item 3): a
step is not a unit of cost, reading a variable or a constant being charged the size of the
value, so a run of $3k$ steps can cost far more than $k$ and finding a value after $3k$
steps says only that the program halted. The clocked universal machine of
Lemma~\ref{lem:universal-tm} decides the same question as a \emph{program}; what is added
here is the Lean-level decision procedure, which is what a computability statement about a
family of verifiers can be built from.

\effortEasy

\begin{lemma}[Semideciders on encoded values]
\label{lem:semidecide-encoded}
\lean{MIPRE.Cost.Data.primrec_decode_prod_nat, MIPRE.Cost.exists_semidecider_of_decode,
  MIPRE.Cost.exists_semidecider_prod_nat}
\leanok
Let $\alpha$ be a type encoded as data with a primitive recursive decoding. Then every
recursively enumerable predicate on $\alpha$ is the halting set of a closed program of the
cost model on the encodings of its arguments. In particular this holds for pairs
$(x, n)$ of a bit string and a number.
\end{lemma}

\begin{proof}
\uses{lem:value-lower-approx}
\leanok
Prefix the bit-string semidecider of Lemma~\ref{lem:value-lower-approx} with the program
that writes the postorder serialization of its input (\texttt{Prog.serProg},
Remark~\ref{rem:compression-abstract}). Serialization and the stack machine
\texttt{Data.parse} invert each other, so the predicate ``decode the parse of this bit
string, and test it'' holds of the serialization of $\mathrm{encode}(a)$ exactly when the
original predicate holds of $a$; and it is recursively enumerable, being the original
predicate precomposed with a computable function.
\end{proof}

\paragraph{Comments.} The detour through a serialization is what keeps the translation into
\texttt{ToPartrec} codes untouched. That translation is delicate for one reason --- a
\texttt{ToPartrec} code cannot see trailing zeros of its input list, while the bit
\texttt{false} encodes as the datum $\mathtt{nil}$ --- and doing it again for each new type
of argument would repeat the delicate step; serializing first does it once.

\effortEasy
